%% Owner: Víctor Martín Castilla
\section{Swarm cluster}
\label{sec:swarm}

All three nodes are Swarm managers, so the cluster keeps working when one
of them fails (\cref{sec:infra-three}). Each node advertises its private
address, so Swarm traffic never leaves the subnet. \Cref{tab:nodes} shows
the nodes as the security review of 17 October found them.

\begin{table}[htbp]
  \centering
  \caption{The Swarm's nodes, as \texttt{docker node ls} reports them.}
  \label{tab:nodes}
  \begin{tabular}{lll}
    \toprule
    Node                     & Status & Manager status \\
    \midrule
    \texttt{swarm-manager-1} & Ready  & Reachable \\
    \texttt{swarm-manager-2} & Ready  & Leader \\
    \texttt{swarm-manager-3} & Ready  & Reachable \\
    \bottomrule
  \end{tabular}
\end{table}

\subsection{Placement}

Pinned services name their node in a placement constraint on its host
name. nginx, for example, runs where
\texttt{node.hostname == swarm-manager-1}. Node labels would read
better, but each one is a manual command on the cluster, while host names
need no step beyond deploying the Compose file.

\subsection{Networks}

Every stack has its own overlay network for internal traffic, such as a
database and its clients. A shared overlay network, \texttt{proxy},
connects nginx with the public services of every stack, which join it
under an alias such as \texttt{app-core} or \texttt{atenea-moodle}. nginx
resolves these names through Docker's DNS on every request, so it starts
even when a stack is missing. \texttt{proxy} is created once, outside any
stack, so that every stack can declare it as external.

nginx publishes ports 80 and 443 through the routing mesh, so every node
listens on them. The security groups only let that traffic reach
\texttt{swarm-manager-1}.

\subsection{Shared disk}

\texttt{swarm-manager-1} runs the kernel NFS server and exports
\texttt{/srv/nfs} to the subnet only. The other two nodes mount it with
NFS~4.2 at the same path, from \texttt{/etc/fstab}, so a bind mount of a
folder under \texttt{/srv/nfs} gives the same files on every node. Swarm
refuses a bind mount whose source does not exist, so a service that
needs a new folder mounts the whole share and creates the folder itself.

The export first squashed root, mapping it to \texttt{nobody} on the
client nodes. The registry, on \texttt{swarm-manager-2}, could not store
images, so the export switched to \texttt{no\_root\_squash}, still
restricted to \texttt{10.0.1.0/24}. The price is that root in any
container that mounts the share acts as root on every shared file.
Services that face the internet should therefore mount only their own
folder, as the Moodle replicas do (\cref{sec:stack-atenea}).

\subsection{Validation}

\begin{itemize}
  \item \textbf{Firewall.} From the internet, only
        \texttt{swarm-manager-1} answers, on ports 22, 80 and 443
        (\cref{sec:infra-network}).
  \item \textbf{Shared disk.} During a redeploy, Swarm moved the app's
        PostgreSQL from \texttt{swarm-manager-1} to
        \texttt{swarm-manager-3}, and the new server found its data
        (\cref{sec:app-nfs}).
  \item \textbf{Overlay network and routing mesh.} The app answers at its
        subdomain through nginx on \texttt{swarm-manager-1}, while its
        replicas run spread over the three nodes (\cref{sec:stack-app}).
\end{itemize}
